\documentclass[10pt,a4paper]{article}

% ---------- encoding & fonts ----------
\usepackage[T1,T2A]{fontenc}
\usepackage[utf8]{inputenc}
% lmodern omitted on purpose: it lacks T2A italic shapes; the default Computer
% Modern (via cm-super) provides full Cyrillic italic and looks identical.
\usepackage[english,russian]{babel}

% ---------- layout ----------
\usepackage[a4paper,top=0.9cm,bottom=0.8cm,left=1.1cm,right=1.1cm]{geometry}
\usepackage{paracol}
\usepackage{titlesec}
\usepackage{enumitem}
\usepackage{xcolor}
\usepackage{hyperref}
\usepackage{fontawesome5}
\usepackage{microtype}

% ---------- colors ----------
\definecolor{accent}{HTML}{2A5DB0}
\definecolor{darktext}{HTML}{1A1A1A}
\color{darktext}

\hypersetup{colorlinks=true,urlcolor=accent,linkcolor=accent}

% ---------- spacing ----------
\setlength{\parindent}{0pt}
\pagestyle{empty}
\setlength{\columnsep}{0.7cm}

% ---------- section headers ----------
\titleformat{\section}{\large\bfseries\color{accent}}{}{0pt}{}
\titlespacing*{\section}{0pt}{5pt}{2pt}

\newcommand{\sidehead}[1]{\vspace{6pt}{\bfseries\color{accent}#1}\par\vspace{2pt}}

% ---------- entry macros ----------
\newcommand{\entry}[3]{%
  \vspace{2pt}%
  \noindent\begin{minipage}[t]{0.73\linewidth}\raggedright{\bfseries #1}\end{minipage}%
  \hfill\begin{minipage}[t]{0.25\linewidth}\raggedleft{\small #2}\end{minipage}\par
  {\itshape #3}\par
  \vspace{0.5pt}%
}
\newcommand{\eduentry}[2]{%
  \vspace{2pt}%
  \noindent\begin{minipage}[t]{0.73\linewidth}\raggedright{\bfseries #1}\end{minipage}%
  \hfill\begin{minipage}[t]{0.25\linewidth}\raggedleft{\small #2}\end{minipage}\par
}

\newlist{tight}{itemize}{1}
\setlist[tight]{leftmargin=12pt,topsep=1pt,itemsep=0.5pt,parsep=0pt,label=\textcolor{accent}{\footnotesize$\bullet$}}

\newcommand{\skillline}[2]{{\bfseries #1}\par{\footnotesize #2}\vspace{3pt}\par}

% кликабельная ссылка обычным тёмным цветом (не синяя)
\newcommand{\plink}[2]{\begingroup\hypersetup{urlcolor=darktext,linkcolor=darktext}\href{#1}{#2}\endgroup}

\begin{document}

\columnratio{0.70}
\begin{paracol}{2}

% ----------------- ЛЕВАЯ КОЛОНКА -----------------
{\fontsize{26}{28}\selectfont Никита Пеганов}\\[2pt]
{\Large\color{accent}Исследователь ИИ}\\[4pt]
{\small Исследователь в области машинного обучения: большие языковые модели, мультиагентные
системы и обучение с подкреплением. Три года промышленной разработки на Python в Яндексе.
Автор нескольких научно-исследовательских работ по нейроморфным вычислениям, RL и LLM-системам.
Сейчас занимаюсь безопасностью GenAI и применением ИИ в психологии.}

\section{Научная работа}

\entry{LLM-трансляция кода, Python\,$\rightarrow$\,Go}{Яндекс, 2026}{Прикладное исследование, ML-инженер}
\begin{tight}
  \item разработал LLM-транслятор, переносящий продакшен-сервисы с Python на Go
  \item спроектировал пайплайн оценки функциональной корректности генерируемого кода
\end{tight}

\entry{\plink{https://github.com/NikPeg/multiagent_llm_strategy}{Мультиагентные LLM-системы}}{2024--2025}{Самостоятельное исследование}
\begin{tight}
  \item сравнил открытые LLM, играющие друг с другом в \plink{https://github.com/NikPeg/llm-graph-mafia-game}{мафию} (локально на A100)
  \item построил мультиагентную LLM-игру для изучения переговорных стратегий
\end{tight}

\entry{\plink{https://github.com/NikPeg/Reinforcement-learning-for-resource-allocation-tasks-in-the-cloud}{RL для распределения ресурсов в облаке}}{Huawei, 2021--2022}{Исследователь, Data Scientist}
\begin{tight}
  \item обучил RL-агента, обрабатывающего 4\,500 запросов/час, всего 1,5 млн запросов
  \item проверил применимость обучения с подкреплением в реальных условиях
\end{tight}

\entry{\plink{https://github.com/NikPeg/synchronization-of-neuromorphic-networks-of-the-close-world-from-the-point-of-view-of-complexes}{Синхронизация нейроморфных сетей тесного мира}}{Школа РАИИ, 2021--2022}{Исследователь}
\begin{tight}
  \item смоделировал нейроморфные сети тесного мира (модель Курамото)
  \item изучил динамику синхронизации при последовательном удалении нейронов
\end{tight}

\section{Проекты}

\entry{\plink{http://212.192.0.8.sslip.io:8080/empathy-ai}{EmpathyAI} --- ИИ-чатбот-психолог}{2025}{Основатель, разработчик}
\begin{tight}
  \item чатбот-психолог с \textbf{100+ активными пользователями}; проактивность и долгосрочная аналитика
  \item 2-й этап конкурса <<Студенческий стартап>> (Фонд содействия инновациям)
\end{tight}

\entry{\plink{https://github.com/NikPeg/GPTScribot}{Scribo} --- LLM-помощник для курсовых}{2024}{Основатель, разработчик}
\begin{tight}
  \item LLM-сервис помогает студентам писать курсовые и дипломы; \textbf{1\,583 пользователя}
\end{tight}

\entry{frAId --- ИИ-супервизор для психотерапевтов}{2023}{Исследователь, разработчик}
\begin{tight}
  \item бот-супервизор анализирует сессии и даёт обратную связь психологам (КПТ)
\end{tight}

\section{Опыт работы}

\entry{Яндекс}{Февраль 2023 --- н.\,в.}{Python-разработчик, инфраструктура}
\begin{tight}
  \item разрабатываю высоконагруженные сервисы и сборщики метрик для L3-управления инфраструктурой
  \item построил модели прогнозирования временных рядов для предсказания нагрузки сервисов
  \item спроектировал аллокатор, распределяющий сервисы между балансерами через решение системы неравенств (cvxpy)
\end{tight}

\section{Образование}

\eduentry{\plink{https://www.hse.ru/ma/ipii/}{НИУ ВШЭ --- магистратура, <<Исследования и предпринимательство в ИИ>>}, средний балл 7,38/10}{2024--2026}
\eduentry{\plink{https://www.hse.ru/ba/se/}{НИУ ВШЭ --- бакалавриат, <<Программная инженерия>>}, средний балл 6,29/10}{2020--2024}

\switchcolumn

% ----------------- ПРАВАЯ КОЛОНКА -----------------
\small
\selectlanguage{russian}
\faMapMarker* ~Москва, Россия\\[2pt]
\faEnvelope ~\href{mailto:peganov.nik@gmail.com}{peganov.nik@gmail.com}\\[2pt]
\faTelegram ~\href{https://t.me/peganov_ns}{t.me/peganov\_ns}\\[2pt]
\faGithub ~\href{https://github.com/NikPeg}{github.com/NikPeg}\\[2pt]
\faLinkedin ~\href{https://linkedin.com/in/nikpeg}{linkedin.com/in/nikpeg}

\sidehead{Научные интересы}
{\footnotesize
Безопасность GenAI (вредоносный код, генерируемый LLM)\,---\,статья в работе\\[3pt]
LLM-агенты и мультиагентные системы\\[3pt]
Обучение с подкреплением\\[3pt]
ИИ в психологии
}

\sidehead{Навыки}
\skillline{Python}{PyTorch, NumPy, pandas, scikit-learn, Jupyter}
\skillline{LLM и NLP}{fine-tuning, RAG, prompt engineering, оценка, локальный inference}
\skillline{Глубокое обучение}{трансформеры, нейросети, генеративные модели}
\skillline{Обучение с подкреплением}{policy optimization, reward design}
\skillline{Мультиагентные системы}{LLM-агенты, моделирование}
\skillline{Прикладной ML}{прогноз временных рядов, выпуклая оптимизация (cvxpy)}
\skillline{Инфраструктура}{Git, Linux, Docker, CUDA / A100}

\sidehead{Языки}
{\footnotesize \textbf{\plink{https://github.com/NikPeg/nikpeg.github.io/blob/main/docs/IELTS.pdf}{English}} --- C1\\[2pt]\textbf{Русский} --- родной}

\sidehead{Достижения}
{\footnotesize
НИР оценены на 9/10 научными руководителями\\[3pt]
Победитель Открытой олимпиады ИТМО по программированию (2020)\\[3pt]
Победитель Межвузовской математической олимпиады (2020)
}

\end{paracol}
\end{document}
